\documentclass[11pt,letterpaper]{article}


\usepackage[breaklinks, pagebackref=true, urlcolor=blue, colorlinks, citecolor=green!50!black, linkcolor=blue,bookmarksopen,bookmarksnumbered,linktocpage=true]{hyperref}

\usepackage[letterpaper, left=1in, right=1in, top=0.9in, bottom=0.9in]{geometry}
\usepackage[utf8]{inputenc}
\usepackage[american]{babel}
\usepackage[normalem]{ulem}
\usepackage{amsmath, amssymb, cases, amsthm}
\usepackage{thmtools}
\usepackage[shortlabels]{enumitem}
\usepackage[framemethod=TikZ]{mdframed}
\usepackage{bbm}
\usepackage{bm}
\usepackage{microtype}
\usepackage{xcolor}
\usepackage{makecell}
\usepackage{mathtools}
\usepackage{float}
\usepackage{modletters}
\usepackage{comment}
\usepackage{multirow}
\usepackage[sort]{cite}
%\usepackage{natbib}
\usepackage[capitalize,noabbrev]{cleveref}
\usepackage{graphics}
\usepackage{xspace}
% pdfsync omitted: editor synchronization only
\usepackage{tocloft}
\setlength\cftbeforesecskip{0.6em}

\usepackage{footnotehyper}
\makesavenoteenv{tabular}

\usepackage{pifont}
\newcommand{\cmark}{\ding{51}}%
\newcommand{\xmark}{\ding{55}}%

\providecommand{\mleftright}{}\let\mleft\left\let\mright\right

\declaretheorem[numberwithin=section,refname={Theorem,Theorems},Refname={Theorem,Theorems}]{theorem}
\declaretheorem[numberwithin=section,refname={Theorem,Theorems},Refname={Theorem,Theorems}]{thm}
\declaretheorem[numberlike=theorem]{lemma}
\declaretheorem[numberlike=theorem]{proposition}
\declaretheorem[numberlike=theorem]{corollary}
\declaretheorem[numberlike=theorem,style=definition]{definition}
\declaretheorem[numberlike=theorem, refname={Claim, Claims}, Refname={Claim, Claims}]{claim}
\declaretheorem[numberlike=theorem,style=remark]{remark}
\declaretheorem[numberlike=theorem,refname={Fact,Facts},Refname={Fact,Facts},name={Fact}]{fact}
\declaretheorem[numberlike=theorem, refname={Question,Questions},Refname={Question,Questions},name={Question}]{question}
\declaretheorem[numberlike=theorem, refname={Problem,Problems},Refname={Problem,Problems},name={Problem}]{problem}
\declaretheorem[numberlike=theorem, refname={Observation,Observations},Refname={Observation,Observations},name={Observation}]{observation}
\declaretheorem[numberlike=theorem, refname={Experiment,Experiments},Refname={Experiment,Experiments},name={Experiment}]{experiment}
\let\primitive\relax
\declaretheorem[numberlike=theorem,refname={Primitive,Primitives},Refname={Primitive,Primitives},name={Primitive}]{primitive}
\declaretheorem[numberlike=theorem,refname={Technique,Techniques},Refname={Technique,Techniques},name={Technique}]{technique}
\declaretheorem[numberlike=theorem,refname={Conjecture,Conjectures},Refname={Conjecture,Conjectures},name={Conjecture}]{conjecture}

\def\cameraready{0}  
\def\final{0}  
\def\iflong{\iffalse}
\ifnum\final=0  

\newcommand{\tawei}[1]{{\color{blue} [{\bf Ta-Wei:} #1]}}
\newcommand{\parth}[1]{{\color{purple} [{\bf Parth:} #1]}}
\newcommand{\TODO}[1]{{\color{blue!50!black} [{\bf Todo:} #1]}}
\else 
\newcommand{\yonggang}[1]{}
\newcommand{\danupon}[1]{}
\newcommand{\sagnik}[1]{}
\newcommand{\todo}[1]{}
\newcommand{\yuval}[1]{}
\newcommand{\jan}[1]{}
\newcommand{\blikstad}[1]{}
\newcommand{\tawei}[1]{}
\newcommand{\parth}[1]{}
\newcommand{\TODO}[1]{}
\fi  

\bibliographystyle{alpha}

\DeclarePairedDelimiter{\ceil}{\lceil}{\rceil}
\newcommand{\eps}{\varepsilon}

\newcommand{\mindeg}{\mathrm{mindeg}}
\newcommand{\nbh}{\mathrm{nbh}}
\newcommand{\spf}{\mathrm{spf}}
\newcommand{\direct}{MDCP }
\newcommand{\best}{\mathsf{best}}
\newcommand{\polylog}{\mathrm{polylog}}
\newcommand{\poly}{\mathrm{poly}}
\newcommand{\outdeg}{\mathrm{outdeg}}
\newcommand{\ones}{\mathrm{ones}}
\DeclareMathOperator*{\var}{Var}
\newcommand{\euler}{\mathrm{e}}
\newcommand{\set}[2][ ]{\{#2 \ifthenelse{\equal{#1}{ }}{ }{~|~#1}\}}
\newcommand{\fvedge}{\mathsf{FindViolatingEdge}}
\newcommand{\XORQ}{\mathsf{XOR}}
\newcommand{\ANDQ}{\mathsf{AND}}

\newcommand{\cc}{\mathrm{cc}}
\newcommand{\Rcal}{\mathcal{R}}
\newcommand{\Pcal}{\mathcal{P}}
\newcommand{\Acal}{\mathcal{A}}
\newcommand{\Scal}{\mathcal{S}}
\newcommand{\Qcal}{\mathcal{Q}}
\newcommand{\pflag}{\mathrm{pflag}}
\newcommand{\R}{\mathbb{R}}
\newcommand{\N}{\mathbb{N}}
\newcommand{\one}{\mathbf{1}}
\newcommand{\inact}{\mathsf{Inactive}}
\newcommand{\act}{\mathsf{Active}}
\newcommand{\cut}{\mathrm{cut}}
\newcommand{\cutd}{\overrightarrow{\mathrm{cut}}}
\newcommand{\Ed}{\overrightarrow{E}}
\newcommand{\mincut}{\lambda}
\newcommand{\tOh}{\widetilde{O}}
\newcommand{\ith}{i^{\scriptsize \mbox{{\rm th}}}}
\newcommand{\jth}{j^{\scriptsize \mbox{{\rm th}}}}



\newcommand{\intpot}[1]{Int_{pot}(#1)}


\newcommand{\ip}[1]{\mleft\langle #1\mright\rangle }
\newcommand{\expt}[2]{\underset{#1}{\mathbb{E}}\mleft[#2\mright]}
\newcommand{\A}{{\mathcal{A}}}
\newcommand{\BPM}{\mathsf{BPM}}
\newcommand{\BMM}{\mathsf{BMM}}
\newcommand{\UBPM}{\mathsf{UBPM}}

\newcommand{\T}{{\mathcal{T}}}

\newcommand{\U}{{\mathcal{U}}}

\newcommand{\bcc}[1]{\textbf{BCAST($#1$)}}
\newcommand{\dw}[1]{#1^{\downarrow}}

\newcommand{\MM}{\mathcal{P}}
\newcommand{\ORQ}{\mathsf{OR}}
\newcommand{\ISQ}{\mathsf{IS}}
\newcommand{\vol}{\mathrm{vol}}
\newcommand{\HPC}{\mathsf{HPC}}
\newcommand{\distance}{\mathsf{dist}}

\newcommand{\rank}{\mathsf{rank}}
\newcommand{\sspan}{\mathsf{span}}
\DeclareMathOperator*{\E}{\mathbb{E}}


\usepackage[most]{tcolorbox}
\tcbuselibrary{skins,breakable}
\tcbset{enhanced jigsaw}
\newcounter{myexample}
\usepackage{xparse}
\usepackage{lipsum}

\makeatletter
\newcommand\footnoteref[1]{\protected@xdef\@thefnmark{\ref{#1}}\@footnotemark}
\makeatother


\renewcommand{\paragraph}[1]{\medskip\noindent{\bf #1}\xspace}






\newcommand{\ground}{U}


\usepackage{tikz}
\usetikzlibrary{arrows}
\usetikzlibrary{arrows.meta}
\usetikzlibrary{shapes}
\usetikzlibrary{backgrounds}
\usetikzlibrary{positioning}
\usetikzlibrary{decorations.markings}
\usetikzlibrary{patterns}
\usetikzlibrary{calc}
\usetikzlibrary{fit}
\usetikzlibrary{decorations}


\renewcommand{\geq}{\geqslant}
\renewcommand{\leq}{\leqslant}
\renewcommand{\le}{\leq}
\renewcommand{\ge}{\geq}
% Parth's macros

\theoremstyle{definition}
\newtheorem{algocf}{Algorithm}
\newenvironment{Algorithm}{\begin{abox}\begin{algocf}}{\end{algocf}\end{abox}}

\Crefname{algocf}{Algorithm}{Algorithms}

\newenvironment{abox}{\begin{tcolorbox}[
		enlarge top by=5pt,
		enlarge bottom by=5pt,
		breakable,
		frame hidden,
		overlay broken = {
			\draw[line width=1pt, black]
			(frame.north west) rectangle (frame.south east);},
		overlay = {
			\draw[line width=1pt, black]
			(frame.north west) rectangle (frame.south east);},
		boxsep=0pt,
		left=4pt,
		right=4pt,
		top=10pt,
		arc=0pt,
		boxrule=1pt,toprule=1pt,
		colback=white
	]}
{\end{tcolorbox}}

\newcommand{\degen}{\textnormal{\textsf{DEGEN}}}

\newcommand{\fail}{\textsc{Reject}}
\newcommand{\true}{\textsc{Accept}}



\providecommand{\given}

\DeclarePairedDelimiter{\pparen}()
\DeclarePairedDelimiter{\card}\vert\vert

\DeclarePairedDelimiterX{\brac}[1][]{
%\renewcommand\given{\nonscript\: \delimsize\vert{} \nonscript\: \mathopen{}}
\ifblank{#1}{\,\cdot\,}{#1}}

\DeclarePairedDelimiterX{\Set}[1]\{\}{
  \renewcommand\given{\nonscript\: \delimsize\vert{} \nonscript\: \mathopen{}}
\ifblank{#1}{\,\cdot\,}{#1}}

\renewcommand{\given}{\nonscript\: \vert \nonscript\:}


\NewDocumentCommand{\Prob}{sO{}E{_}{{}}m}{%
    \Pr_{#3}
    \IfBooleanTF{#1}
    {\brac*{#4}}
    {\brac[#2]{#4}}
}

\DeclarePairedDelimiterXPP{\Ot}[1]{\widetilde{O}}(){}{#1}
\DeclarePairedDelimiterXPP{\Omgt}[1]{\widetilde{\Omega}}(){}{#1}

\DeclarePairedDelimiterXPP{\Var}[1]{\var}[]{}{#1}

\newcommand{\dsi}{\cD_\SI}
\newcommand{\pisi}{\pi_\SI}

\newcommand{\ptd}{\Lambda}
\DeclarePairedDelimiterXPP{\PTD}[2]{\Lambda}(){}{#1, #2}

\DeclareMathOperator*{\distrib}{dist}
\NewDocumentCommand{\Dist}{sO{}E{_}{{}}m}{%
  \mu_{#3}
  \IfBooleanTF{#1}
  {\pparen*{#4}}
  {\pparen[#2]{#4}}
}

\DeclareMathOperator*{\expect}{\mathbb{E}}
\NewDocumentCommand{\Exp}{sO{}E{_}{{}}m}{%
  \expect_{#3}
  \IfBooleanTF{#1}
  {\brac*{#4}}
  {\brac[#2]{#4}}
}

\newcommand{\re}{\rv{e}}

\newcommand{\estar}{e^{\ast}}
\newcommand{\restar}{\re^{\ast}}

\newcommand{\rPi}{\rv{\Pi}}

\newcommand{\Perm}{\Gamma}
\newcommand{\rPerm}{\rv{\Gamma}}

\newcommand{\uni}{\mathcal{U}}

\newcommand{\qpi}{q^{\Pi}}

\newcommand{\score}{\mathsf{Score}}
\newcommand{\totscore}{\mathsf{TScore}}


\DeclareMathOperator*{\inform}{\mathbb{I}}
\DeclarePairedDelimiterXPP{\Inf}[1]{\inform}(){}{#1}

\DeclareMathOperator*{\entropy}{\mathbb{H}}
\DeclarePairedDelimiterXPP{\Ent}[1]{\entropy}(){}{#1}

\DeclareMathOperator{\myIC}{\mathsf{IC}}
\DeclareMathOperator{\myCC}{\mathsf{CC}}
\newcommand{\ExtIC}{\myIC^{\textnormal{ext}}}
\newcommand{\IntIC}{\myIC^{\textnormal{int}}}

\newcommand{\rPTD}{\mathbf{\Lambda}}

\newcommand{\Leq}[1]{{\underset{\textnormal{#1}}\leq}}
\newcommand{\Geq}[1]{{\underset{\textnormal{#1}}\geq}}
\newcommand{\Eq}[1]{{\underset{\textnormal{#1}}=}}



\newcommand{\bhpc}{\ensuremath{\textnormal{\textsf{BHPC}}}\xspace}
\newcommand{\bmhpc}{\ensuremath{\textnormal{\textsf{BMHPC}}}\xspace}
\newcommand{\hpc}{\ensuremath{\textnormal{\textsf{HPC}}}\xspace}
\newcommand{\mhpc}{\ensuremath{\textnormal{\textsf{MHPC}}}\xspace}
\newcommand{\setint}{\ensuremath{\textnormal{\textsf{SetInt}}}\xspace}
\newcommand{\defc}{\mathtt{def}}
\newcommand{\odeg}{\text{odeg}}
\newcommand{\Vaux}{V_{\text{aux}}}
\newcommand{\Eaux}{E_{\text{aux}}}
\newcommand{\prantar}[1]{{\color{blue} [{\bf Prantar:} #1]}}
\newcommand{\mypar}[1]{\smallskip\noindent{\sffamily\bfseries #1.}~}
\newcommand{\floor}[1]{\lfloor #1 \rfloor}




\newcommand{\cmax}{\ensuremath{\textnormal{\textsf{c}}_{\textnormal{rs}}}}
\newcommand{\protMO}{\ensuremath{\Prot_{\mo}}\xspace}
\newcommand{\bProtMO}{\ensuremath{\bm{\Prot}_{\mo}}\xspace}
\newcommand{\protBHH}{\ensuremath{\Prot_{\bhh}}\xspace}
\newcommand{\bProtBHH}{\ensuremath{\bm{\Prot}_{\bhh}}\xspace}


\newcommand{\Paren}[1]{\Big(#1\Big)}
\newcommand{\Bracket}[1]{\Big[#1\Big]}
\newcommand{\bracket}[1]{\mleft[#1\mright]}
\newcommand{\paren}[1]{\ensuremath{\mleft(#1\mright)}\xspace}

\newenvironment{tbox}{\begin{tcolorbox}[
		enlarge top by=5pt,
		enlarge bottom by=5pt,
		 breakable,
		 boxsep=0pt,
                  left=4pt,
                  right=4pt,
                  top=10pt,
                  boxrule=1pt,toprule=1pt,
                  colback=white,
                  arc=-1pt,
        %          drop shadow={black,opacity=1},
                 % arc=0pt
                  ]%%
	}
{\end{tcolorbox}}



\newcommand{\Prot}{\ensuremath{\Pi}}
\newcommand{\prot}{\ensuremath{\pi}}
\newcommand{\Rpub}{{R}^{pub}}
\newcommand{\Rpri}{{R}^{pri}}

\newcommand{\bProt}{\bm{\Prot}}
\newcommand{\bRpub}{\bm{\Rpub}}
\newcommand{\bRpri}{\bm{\Rpri}}


\newcommand{\bV}{\bm{V}}
\newcommand{\bX}{\bm{X}}
\newcommand{\bY}{\bm{Y}}
\newcommand{\bZ}{\bm{Z}}
\newcommand{\bA}{\bm{A}}
\newcommand{\bP}{\bm{P}}
\newcommand{\bR}{\bm{R}}
\newcommand{\bM}{\bm{M}}
\newcommand{\bI}{\bm{I}}
\newcommand{\bJ}{\bm{J}}
\newcommand{\bSigma}{\bm{\sigma}}
\newcommand{\bG}{\bm{G}}

\newcommand{\bB}{\ensuremath{\bm{B}}}
\newcommand{\bC}{\ensuremath{\bm{C}}}
\newcommand{\bD}{\ensuremath{\bm{D}}}

\newcommand{\supp}[1]{\ensuremath{\textsc{supp}(#1)}}




\newtheorem{mdresult}{Result}
\newenvironment{result}{\begin{mdframed}[backgroundcolor=lightgray!40,topline=false,rightline=false,leftline=false,bottomline=false,innertopmargin=2pt]\begin{mdresult}}{\end{mdresult}\end{mdframed}}


%-------------------------- Textbox ---------------------------------

\newcommand{\textbox}[2]{
{
\begin{tbox}
\textbf{#1}
{#2}
\end{tbox}
}
}

\newcommand{\rA}{\rv{A}}
\newcommand{\rB}{\rv{B}}
\newcommand{\rC}{\rv{C}}
\newcommand{\rD}{\rv{D}}

\newcommand{\rX}{\rv{X}}
\newcommand{\rY}{\rv{Y}}
\newcommand{\rI}{\rv{I}}

\newcommand{\rv}[1]{\ensuremath{\mathsf{#1}}}

\newcommand{\II}{\ensuremath{\mathbb{I}}}
\newcommand{\HH}{\ensuremath{\mathbb{H}}}
\newcommand{\mi}[2]{\ensuremath{\II(#1 \,; #2)}}
\newcommand{\en}[1]{\ensuremath{\HH(#1)}}

\newcommand{\itfacts}[1]{Fact~\ref{fact:it-facts}-(\ref{part:#1})\xspace}

\newcommand{\DD}{\ensuremath{\mathbb{D}}}

\newcommand{\distribution}[1]{\ensuremath{\textnormal{dist}(#1)}\xspace}

\newcommand{\XX}{\ensuremath{\mathcal{X}}}
\newcommand{\YY}{\ensuremath{\mathcal{Y}}}


\newcommand{\HPCk}{\ensuremath{\HPC_k}\xspace}

\newcommand{\protHPC}{\ensuremath{\prot_{\HPC}}}

\newcommand{\distHPC}{\ensuremath{\dist_{\HPC}}}

\newcommand{\unif}{\ensuremath{\mathcal{U}}}

\newcommand{\rZ}{\rv{Z}}
\newcommand{\rt}{\rv{t}}
\newcommand{\rT}{\rv{T}}
\newcommand{\rE}{\rv{E}}
\newcommand{\rP}{\rv{P}}
\newcommand{\rProt}{\rv{\Prot}}

\newcommand{\distSI}{\ensuremath{\dist}_{\ensuremath{\textnormal{\textsf{SI}}}}\xspace}


\newcommand{\protSI}{\ensuremath{\prot}_{\ensuremath{\textnormal{\textsf{SI}}}}\xspace}
\newcommand{\ProtSI}{\ensuremath{\Prot}_{\ensuremath{\textnormal{\textsf{SI}}}}\xspace}
\newcommand{\rProtSI}{\ensuremath{\rv{\Prot}}_{\ensuremath{\textnormal{\textsf{SI}}}}\xspace}

\newcommand{\distPI}{\ensuremath{\dist}_{\ensuremath{\textnormal{\textsf{PI}}}}\xspace}
\newcommand{\protPI}{\ensuremath{\prot}_{\ensuremath{\textnormal{\textsf{PI}}}}\xspace}
\newcommand{\ProtPI}{\ensuremath{\Prot}_{\ensuremath{\textnormal{\textsf{PI}}}}\xspace}
\newcommand{\rProtPI}{\ensuremath{\rv{\Prot}}_{\ensuremath{\textnormal{\textsf{PI}}}}\xspace}


\newcommand{\SI}{\ensuremath{\textnormal{\textsf{Set-Int}}}\xspace}
\newcommand{\PI}{\ensuremath{\textnormal{\textsf{Pair-Int}}}\xspace}


\newcommand{\errs}{\ensuremath{\textnormal{errs}}\xspace}




\newcommand{\rK}{\rv{K}}
\newcommand{\rJ}{\rv{J}}
\newcommand{\rL}{\rv{L}}
\newcommand{\rS}{\rv{S}}
\newcommand{\rM}{\rv{M}}
\newcommand{\rR}{\rv{R}}
\newcommand{\rbarS}{\rv{\barS}}



\newcommand{\distDisjN}{\ensuremath{\distDisj^{\textnormal{\textsf{N}}}}\xspace}
\newcommand{\distDisjY}{\ensuremath{\distDisj^{\textnormal{\textsf{Y}}}}\xspace}

\newcommand{\mii}[3]{\ensuremath{\II_{#3}(#1 \,; #2)}}

\newcommand{\disjin}[4]{ \ensuremath{\begin{bmatrix} #1&\hspace{-8pt}#2 \\ #3&\hspace{-10pt}#4 \end{bmatrix}}}
\renewcommand{\disjin}[4]{\textnormal{[\ensuremath{#1#2,#3#4}]}}

\newcommand{\barS}{\overline{S}}

\newcommand{\hd}[2]{\ensuremath{\textnormal{h}(#1,#2)}\xspace}
\newcommand{\hdt}[2]{\ensuremath{\textnormal{h}^2(#1,#2)}\xspace}


\newcommand{\CC}[2]{\ensuremath{\textnormal{\textsf{CC}}_{#2}(#1)}\xspace}
\newcommand{\ICost}[2]{\ensuremath{\textnormal{\textsf{IC}}_{#2}(#1)}\xspace}
\newcommand{\IC}[2]{\ICost{#1}{#2}}


\newcommand{\rTheta}{\rv{\Theta}}



\newcommand{\comp}{\succ_{\ProtSI}}
\newcommand{\compr}{\prec_{\ProtSI}}
\newcommand{\factor}{3}
%%%%
